% Draft: §3, Det--Head, for the restructured paper (rewrite method agreed
% 17 Oct 2026; see the evidence repo's DECISIONS.md). Drafted from the old
% draft's §§2, 3.1--3.3, 5, 7 and 8, CGEL ch. 5 §9, and Payne et al. (2007);
% revised after the muse review (reviews/muse-sec3-20261017-135000/TRIAGE.md).
% Verified detail is to be checked in a later pass, as Brett decided.
%
% Labels it relies on: sec:analyses (§2), sec:mod-head (§4), sec:once (§5),
% sec:fused-relatives (§7). Not yet \input into the main file.

\section{Determinatives standing alone: Det--Head}\label{sec:det-head}

A determinative can make up a whole noun phrase, or head one that also contains a complement or modifier, and so can an independent genitive such as \mention{mine} (e):

\ea\label{ex:dh}
\ea \mention{I'll take \bk{some}.}
\ex \mention{\bk{Few} would dare underestimate Aaron Finch's men.}
\ex \mention{some of the wine}, \mention{many of them}, \mention{whichever of them}
\ex \mention{hardly anyone present}, \mention{something reliable}
\ex \mention{\bk{Mine} are ready.}
\z\z

For \textit{CGEL} each of these has a fused Det--Head: one phrase functions at once as determiner of the NP and as head of its nominal \citep[410--423]{huddleston2002}. On ordinary Head it's a pure head, with no second function (Section~\ref{sec:analyses}). The question here is whether anything in the form of these expressions shows the determiner function.

Some things aren't at issue. Both analyses take these to be complete NPs with no silent noun. They have the functions of NPs (\mention{Some left}, \mention{I see some}, \mention{with some}), there's no noun to restore (\ungram{\mention{many ones of us}}; pp.~420--421), and often there's none in the context either.\footnote{Of 94 NOW lines for \mention{few would}, 82 have no noun phrase or named set in the displayed context that could supply the set, as in (\ref{ex:dh}b) (\textit{Daily Times}, Pakistan, 20 October 2022); 3 have one, 2 probably do, and 7 can't be decided from the window.} Nor is the lack of a further determiner evidence for the determiner function. \textit{CGEL} gives \mention{She left}, \mention{Kim left}, and \mention{People left} NPs with no determiner, and ordinary Head gives \mention{Some left} the same structure.

Partitives add nothing here. In \mention{some of the wine}, the \mention{of}-phrase is a complement inside the nominal \citep[411--412]{huddleston2002}, so \mention{some} is the ultimate head on either analysis. The claim about bare uses extends to them unchanged, including \mention{whichever of them} in the fused relative of Section~\ref{sec:fused-relatives}. So do relative clauses.\footnote{\textcite[184--185]{wu2021exocentric} finds a dilemma for fusion in \mention{For \bk{some who made the journey into Wales}, the experience was profound} (BNC): fusion makes \mention{some}, a determiner, the antecedent of the relative clause, against the principle that determiners don't count as antecedents, and he adds an empty noun to resolve it. On ordinary Head, \mention{some} heads the nominal and the relative clause modifies that nominal, with nothing added.}

\textit{CGEL} and \textcite{Payne2007} point to three places where the determiner function might show: the correspondence with the determiner use; the forms of \mention{none}, \mention{mine}, and the compounds; and case. None of them shows it. One fact of distribution does favour fusion, and it comes with the compounds.

First, the correspondence. \textit{CGEL} keeps the determiner function because the word \enquote{can occur in the dependent position in NPs where it is not fused with the head} (p.~410). It keeps the category because the word takes the same dependents in both uses (\mention{not much cheese}, \mention{not much of the cheese}; pp.~420--421). Ordinary Head keeps the category, and with it both facts: \mention{much} is the same determinative in both uses, and it takes the dependents a determinative takes.

What the correspondence adds beyond the category is a list of the determinatives that can stand alone. \mention{Some}, \mention{each}, and \mention{either} can, and \mention{every} normally can't (\ungram{\mention{every of the students}}), though NOW has two published lines with it standing alone, whose status I leave open (\mention{Every has a story}, \textit{Patch}, 2025; \mention{Every has different blocking mechanics}, \textit{Guardian}, 2024); \mention{no} does so only in the form \mention{none}. \textcite[581]{Payne2007} state the list as a \enquote{lexical distinction between determinatives that allow [fusion of functions] and those that do not}. Ordinary Head states the same list as a permission to head a nominal with no determiner. Neither analysis gets it free.

\textit{CGEL} considers ordinary Head too, as \enquote{change of function without change of category}, and rejects it because \mention{nylon} in \mention{I prefer cotton shirts to \bk{nylon}} is \enquote{understood as having modifier function} (p.~421). Section~\ref{sec:analyses} answers that objection: an interpretation alone doesn't establish a syntactic function. For determinatives the objection has a further gap. It contrasts a word's use as a pure head with its fused use, and \textit{CGEL} gives determinatives no use as a pure head.

Second, form. Payne et al.\ find formal differences between fused and non-fused functions here, as in fused relatives: \mention{no} \enquote{has a specialized morphological variant \mention{none} which is restricted to the fused function} (p.~581). A difference of form doesn't say which function it marks, though, and the genitive pronouns supply the test. They alternate the same way (\mention{my}, \mention{mine}), and \textit{CGEL} states the conditioning for both pairs at once: \enquote{the dependent forms are used in (pure) determiner function, the independent ones in head function} \citep[410]{huddleston2002}.

\mention{Mine} occurs with a determiner function and without one. The independent genitive pronouns \enquote{occur in fused heads and also pure heads} (p.~410), pure heads being heads with no second function, and predicative \mention{mine} in \mention{Everything that is \bk{mine} is yours} has \enquote{no basis for a fused-head analysis} (p.~469). Its form is the same, so the form marks the head function. \mention{None} has no pure-head use to test (p.~410), so its form is assigned to the same conditioning on the model of \mention{mine}, not shown independently. Ordinary Head then states both alternations as \textit{CGEL} does, by function: Det against Head.

Compounds raise the question of form again. \textit{CGEL} analyses \mention{anyone}, \mention{something}, and the rest as fused determiner-heads \citep[423]{huddleston2002}. Payne et al.\ take them to be determinatives because \enquote{they take exactly the same pre-head modifiers as the determinative bases they contain} (p.~581): \mention{not everything} beside \mention{not every house}, \mention{hardly anyone present} beside \mention{hardly any writer present}. That fixes the category, which ordinary Head keeps.

On ordinary Head, \mention{not} and \mention{hardly} are modifiers of the whole NP, the kind \textit{CGEL} calls peripheral \citep[430--431]{huddleston2002}, and the compound's base still decides which ones it takes. Payne et al.\ also take the compounds' morphology, a determinative base with a nominal base, to mirror fusion (pp.~565, 587--588). What the two bases mirror is the compound's meaning, a quantifier with a sort (\enquote*{person}, \enquote*{thing}), and both analyses give it that meaning. Simple \mention{some} and \mention{few} have no nominal base, and for Payne et al.\ they're fused just the same.\footnote{Payne et al.'s own note treats the nominal base of \mention{once} as a reanalysis of a genitive ending (p.~588 n.~30); see Section~\ref{sec:once}.}

The compounds' modifiers divide by side: adverbs before, adjectives after (\mention{hardly anyone present}, \mention{something reliable}). The adjective is a \term{restrictor}, a single post-head dependent that doesn't recur and can be a simple adjective, of the kind that would normally precede a noun (p.~423; \citealp[583]{Payne2007}). Payne et al.\ count fusion's account of the division \enquote{a particular strength} (p.~583). On ordinary Head the compound, as head of the nominal, takes the restrictor inside the nominal, and the grouping differs: ordinary Head gives \mention{\bk{hardly \bk{anyone present}}}, fusion \mention{\bk{\bk{hardly anyone} present}}.

\textit{CGEL}'s peripheral modifiers can scope over a whole NP, as in \mention{After \bk{hardly a moment's hesitation} he agreed} \citep[437]{huddleston2002}, so a fusion grammar can assign the first grouping too, and the grouping doesn't decide between them.

Nor does fusion account for the restrictor's position by itself. \textit{CGEL} puts the adjective after the head \enquote{because determiner and head are fused}, which rules out pre-head dependents (p.~423). But the restrictor stays where nothing is fused. \textit{CGEL} treats \mention{something} in \mention{a little something} as a noun formed by conversion (p.~423 n.~43), and a single adjective follows it (\mention{a little something extra}), as one does in \mention{a pre-emptive something fabulous} (CBS, 2007). With a common noun as head, a postpositive adjective occurs only \enquote{under restricted conditions} (p.~529).\footnote{In COCA (18 October 2026), 46 of 50 lines with an article, an adjective, \mention{somebody} or \mention{something}, and a second adjective have it after \mention{something}: 43 \mention{a little something} (\mention{extra} in most), plus \mention{a certain something special} (cafnrnews.com, 2012), \mention{a pre-emptive something fabulous}, and \mention{a special something unique to \ldots} (1996), whose adjective phrase has a complement and would follow any head (p.~445). The evidence beyond the formula is thin, and none has a personal compound. \mention{Else} survives in the same frame (\mention{a certain someone else}, \textit{Community}, 2010), but \textit{CGEL} doesn't derive it from fusion, since \mention{all}, \mention{much}, and \mention{little} take it too (p.~423).} So the position goes with the compound, whatever its functions. A defender of fusion can let the converted noun inherit the compound's restrictor, but that's a statement about the word, of the kind ordinary Head makes for every use. The division gives fusion no advantage, then, though it gives ordinary Head none either.

The other half of the division does favour fusion, and it's the one trace of the determiner function in the distribution of these words. Bare determinatives take no pre-head adjective (\ungram{\mention{nice anyone}}, \ungram{\mention{Lucky few escaped}}), apart from a few exclamations such as \mention{Poor everybody!} (Section~\ref{sec:mod-head}). With the determiner and head functions fused, an adjective would have to come before the determiner, so fusion derives the gap from adjacency, where ordinary Head has to state it, as an analysis with a silent head does \citep[584]{Payne2007}.

The adjective is fine once a determiner precedes (\mention{the lucky few}). It's also fine in \mention{a lucky few}, which \textit{CGEL}'s analyses of \mention{a few} allow two ways \citep[353, 392]{huddleston2002}. Section~\ref{sec:mod-head} shows that on either way the adjective doesn't come before a bare determinative, so the gap stands.

Third, case, where \textit{CGEL} itself says the determiner function would show. In \mention{many people's expenses}, the genitive \mention{'s} attaches to \mention{people}, the head of \mention{many people's}, and not to its determiner, \mention{many}. On the fused analysis a bare \mention{many} is both head and determiner, so \ungram{\mention{many's expenses}} would put the genitive on a determiner. \textit{CGEL} reports that very gap, calling \mention{all's}, \mention{many's}, \mention{few's}, and \mention{some's} \enquote{completely excluded}. Its explanation is fusion: \enquote{the form is simultaneously determiner and head, and for full acceptability its case should match both functions} \citep[420]{huddleston2002}.

\textit{CGEL} offers this against analysing these words as pronouns, an analysis that gives no reason for the gap (p.~420). Ordinary Head gives no reason either, so the argument bears on it too, and it's the most direct of the three.

For \mention{all}, \mention{many}, \mention{few}, and \mention{some}, the gap is all but complete in usage,\footnote{COCA searches for each word followed by \mention{'s}, one search per word, give the following genitives on the word standing alone, out of the lines each search returned: \mention{either} 17 of 30 lines, \mention{neither} 6 of 23, \mention{each} 12 of 12, \mention{many} 3 of 76, \mention{some} 1 of 24, \mention{few} none of 22, and \mention{all} none in a display-capped sample of 499 lines (nine genitives there are on \mention{you all}). The other lines are mostly contractions of \mention{is} or \mention{has} (\mention{Either's fine}, \mention{many's the time}) and names. The counts aren't adjusted for how often each word stands alone. Thirteen of the 35 genitives of \mention{either}, \mention{neither}, and \mention{each} come from edited prose (fiction, magazines, newspapers, academic writing). The four of \mention{many} and \mention{some} all come from the web; two seem fine to me (\mention{to \bk{many's} suprise \textup{[sic]} he won}, \mention{out of \bk{so many's} reach}), the other two questionable.} but it doesn't hold across the fused determiner-heads. \mention{Either's}, \mention{neither's}, and \mention{each's} occur in edited prose, and I find them acceptable: \mention{won't substantially increase \bk{either's} share of its market} (\textit{TIME}, 2000), \mention{\bk{Neither's} numbers are all that encouraging} (\textit{Washington Post}, 2014), \mention{the highlight of \bk{each's} day} (\textit{Atlantic}, 1995). \textit{CGEL} calls \mention{each's} \enquote{occasionally attested but of somewhat questionable status} (p.~420). COCA's twelve genitives of \mention{each} fit the first half of that, but five of them are in edited prose, which fits the second less well.

\textit{CGEL}'s own examples of genitive marking on the head include \mention{everyone's} and \mention{somebody's} (p.~479), compounds it analyses as fused determiner-heads (p.~423). A defender of fusion could put their genitive down to the nominal base, \mention{-one} or \mention{-body}.

So the genitive divides these words, and the argument from case faces a dilemma. If \mention{many's} is out because its case would have to match two functions (p.~420), then \mention{either's}, \mention{neither's}, and \mention{each's} should be out too, since \textit{CGEL} gives those words the same two functions. If something else keeps \mention{many's} out, the gap is no evidence for the determiner function. Either way, the two functions don't predict where the division falls, though that doesn't show that function plays no part. As stated, neither analysis predicts it: fusion has to exempt \mention{either}, \mention{neither}, \mention{each}, and the compounds unless their nominal base carries the genitive, and ordinary Head has to list the words that don't take it, as it lists the reflexive pronouns, which have no genitive either (p.~426).

The division may follow meaning. The words that take the genitive, the compounds included, are singular, and those that don't are words of quantity (\mention{all}, \mention{many}, \mention{few}, \mention{some}); if that's what decides, neither analysis needs a list. Either way, what the form varies with is the word, not its functions. A comparison that could show a function would keep the word and change its functions, and for these words there's none to make, since \textit{CGEL} gives them no use as a pure head.

Section~\ref{sec:fused-relatives} finds that kind of evidence. After a preposition, \mention{whomever} is chosen over \mention{whoever} far more often when it isn't the subject of the relative clause than when it is. The word and its function in the main clause stay the same, and the form follows the function inside the relative clause. On the diagnostic \textit{CGEL} itself proposes, then, the determiner function doesn't show in fused determiner-heads, and the second function in fused relatives does.

Both analyses, then, list the same words and forms. The trade is a constituent with two functions against one stated condition on pre-head adjectives, of a kind the grammar already states for pronouns (Section~\ref{sec:mod-head}). Apart from the pre-head adjective, nothing in the form of these expressions shows a determiner function. That fact is a real advantage for fusion, and Section~\ref{sec:costs} weighs it against what ordinary Head gains.
